% GENERATED FILE. Edit canonical lessons, then run npm run generate:books.
% canonical-source-hash: fnv1a64:009cdf5931503512
% canonical-lessons: TE-C470-kalusyam, TE-C470-vyarthalu, TE-C470-vayuvu, TE-C470-kalusitam-ceyu, TE-C470-svacchamaina

\chapter{Dirty Air and Water}
\label{ch:te-dirty-air-and-water}
\hlchaptermodality{470}

\begin{chapteropening}
\textbf{By the end of this chapter:} \emph{I can say pollution, waste, a gas, to pollute and pure.}

Complete the last lesson of chapter 470: I can say pollution, waste, a gas, to pollute and pure.
\end{chapteropening}

\section[\texorpdfstring{\te{కాలుష్యం} (kāluṣyaṁ)}{కాలుష్యం (kāluṣyaṁ)}]{\te{కాలుష్యం} (kāluṣyaṁ) — pollution}

\label{lesson:TE-C470-kalusyam}



\begin{quote}
Before the new one: say the Telugu for a rush, then the Telugu for a mix-up.
\end{quote}

\subsection*{You'll want to know: \te{కాలుష్యం}}
\textbf{\te{కాలుష్యం}} — \emph{kāluṣyaṁ} — \textquotedblleft{}pollution\textquotedblright{}.

\begin{quote}
\te{పెద్ద} \te{నగరాల్లో} \te{గాలి} \te{కాలుష్యం} \te{ఎక్కువగా} \te{ఉంది}.
\end{quote}

\begin{quote}
\emph{pedda nagarāllō gāli kāluṣyaṁ ekkuvagā undi.}
\end{quote}

\begin{quote}
Air pollution is high in the big cities.
\end{quote}

\te{గాలి} \te{కాలుష్యం}, air pollution; \te{నీటి} \te{కాలుష్యం}, water pollution.

\begin{cousinweb}
\textbf{\te{కాలుష్యం}} (\emph{kāluṣyaṁ}) is borrowed from Sanskrit \emph{kāluṣya}, \textquotedblleft{}foulness, dirtiness, turbidness\textquotedblright{} (Monier-Williams), the noun of \emph{kaluṣa}, \textquotedblleft{}foul\textquotedblright{}.
\end{cousinweb}

\begin{grammarlens}[title={What you've built}]
One more word for getting around.
\end{grammarlens}

\subsection*{Your turn}
\begin{itemize}
\raggedright
  \item \emph{Say it:} \emph{kāluṣyaṁ}
  \item \emph{Say it:} \emph{kāluṣyaṁ}, once more
  \item \emph{Say it:} \emph{pedda nagarāllō gāli kāluṣyaṁ ekkuvagā undi.} — Air pollution is high in the big cities.
  \item \emph{Recall:} say the Telugu for a rush, then the Telugu for a mix-up, then say \emph{kāluṣyaṁ} again
\end{itemize}

\begin{tcolorbox}[breakable,colback=teal!4,colframe=teal!35!black,title={Before you move on}]
What does \te{కాలుష్యం} mean? (\textquotedblleft{}Pollution\textquotedblright{}.) Say it once more.
\end{tcolorbox}
\section[\texorpdfstring{\te{వ్యర్థాలు} (vyarthālu)}{వ్యర్థాలు (vyarthālu)}]{\te{వ్యర్థాలు} (vyarthālu) — waste}

\label{lesson:TE-C470-vyarthalu}



\begin{quote}
Before the new one: say the Telugu for a mix-up, then the Telugu for pollution.
\end{quote}

\subsection*{You'll want to know: \te{వ్యర్థాలు}}
\textbf{\te{వ్యర్థాలు}} — \emph{vyarthālu} — \textquotedblleft{}waste\textquotedblright{}.

\begin{quote}
\te{కర్మాగారం} \te{వ్యర్థాలను} \te{నదిలో} \te{వదులుతోంది}.
\end{quote}

\begin{quote}
\emph{karmāgāraṁ vyarthālanu nadilō vadulutōndi.}
\end{quote}

\begin{quote}
The factory is letting its waste into the river.
\end{quote}

Mostly plural: \te{ప్లాస్టిక్} \te{వ్యర్థాలు}, plastic waste. \te{వ్యర్థం} alone means useless, in vain.

\begin{cousinweb}
\textbf{\te{వ్యర్థాలు}} (\emph{vyarthālu}) is borrowed from Sanskrit \emph{vyartha}, \textquotedblleft{}useless, without purpose\textquotedblright{} (Monier-Williams), from \emph{vi-}, \textquotedblleft{}without\textquotedblright{}, and \emph{artha}, \textquotedblleft{}use, purpose\textquotedblright{}; useless things became waste.
\end{cousinweb}

\begin{grammarlens}[title={What you've built}]
One more word for getting around.
\end{grammarlens}

\subsection*{Your turn}
\begin{itemize}
\raggedright
  \item \emph{Say it:} \emph{vyarthālu}
  \item \emph{Say it:} \emph{vyarthālu}, once more
  \item \emph{Say it:} \emph{karmāgāraṁ vyarthālanu nadilō vadulutōndi.} — The factory is letting its waste into the river.
  \item \emph{Recall:} say the Telugu for a mix-up, then the Telugu for pollution, then say \emph{vyarthālu} again
\end{itemize}

\begin{tcolorbox}[breakable,colback=teal!4,colframe=teal!35!black,title={Before you move on}]
What does \te{వ్యర్థాలు} mean? (\textquotedblleft{}Waste\textquotedblright{}.) Say it once more.
\end{tcolorbox}
\section[\texorpdfstring{\te{వాయువు} (vāyuvu)}{వాయువు (vāyuvu)}]{\te{వాయువు} (vāyuvu) — a gas}

\label{lesson:TE-C470-vayuvu}



\begin{quote}
Before the new one: say the Telugu for pollution, then the Telugu for waste.
\end{quote}

\subsection*{You'll want to know: \te{వాయువు}}
\textbf{\te{వాయువు}} — \emph{vāyuvu} — \textquotedblleft{}a gas\textquotedblright{}.

\begin{quote}
\te{ఆ} \te{కర్మాగారం} \te{నుంచి} \te{విష} \te{వాయువు} \te{బయటకు} \te{వచ్చింది}.
\end{quote}

\begin{quote}
\emph{ā karmāgāraṁ nuñci viṣa vāyuvu bayaṭaku vaccindi.}
\end{quote}

\begin{quote}
A poisonous gas came out of that factory.
\end{quote}

A formal word for a gas; in Sanskrit it is the wind. In a kitchen people say \te{గ్యాస్}.

\begin{cousinweb}
\textbf{\te{వాయువు}} (\emph{vāyuvu}) is borrowed from Sanskrit \emph{vāyu}, \textquotedblleft{}wind, air\textquotedblright{} (CDIAL 11544); modern Telugu uses it for a gas.
\end{cousinweb}

\begin{grammarlens}[title={What you've built}]
One more word for getting around.
\end{grammarlens}

\subsection*{Your turn}
\begin{itemize}
\raggedright
  \item \emph{Say it:} \emph{vāyuvu}
  \item \emph{Say it:} \emph{vāyuvu}, once more
  \item \emph{Say it:} \emph{ā karmāgāraṁ nuñci viṣa vāyuvu bayaṭaku vaccindi.} — A poisonous gas came out of that factory.
  \item \emph{Recall:} say the Telugu for pollution, then the Telugu for waste, then say \emph{vāyuvu} again
\end{itemize}

\begin{tcolorbox}[breakable,colback=teal!4,colframe=teal!35!black,title={Before you move on}]
What does \te{వాయువు} mean? (\textquotedblleft{}A gas\textquotedblright{}.) Say it once more.
\end{tcolorbox}
\section[\texorpdfstring{\te{కలుషితం} \te{చేయు} (kaluṣitaṁ cēyu)}{కలుషితం చేయు (kaluṣitaṁ cēyu)}]{\te{కలుషితం} \te{చేయు} (kaluṣitaṁ cēyu) — to pollute}

\label{lesson:TE-C470-kalusitam-ceyu}



\begin{quote}
Before the new one: say the Telugu for waste, then the Telugu for a gas.
\end{quote}

\subsection*{You'll want to know: \te{కలుషితం} \te{చేయు}}
\textbf{\te{కలుషితం} \te{చేయు}} — \emph{kaluṣitaṁ cēyu} — \textquotedblleft{}to pollute\textquotedblright{}.

\begin{quote}
\te{కర్మాగారాలు} \te{నది} \te{నీటిని} \te{కలుషితం} \te{చేశాయి}.
\end{quote}

\begin{quote}
\emph{karmāgārālu nadi nīṭini kaluṣitaṁ cēśāyi.}
\end{quote}

\begin{quote}
The factories polluted the river water.
\end{quote}

Past (I …): \te{కలుషితం} \te{చేశాను}, \emph{kaluṣitaṁ cēśānu}.

\te{కలుషితం}, polluted, + \te{చేయు}. \te{కలుషితం} \te{అవు}, to become polluted.

\begin{cousinweb}
\textbf{\te{కలుషితం} \te{చేయు}} (\emph{kaluṣitaṁ cēyu}) is a Telugu phrase: \te{కలుషితం}, borrowed from Sanskrit \emph{kaluṣita}, \textquotedblleft{}made foul, impure\textquotedblright{} (Monier-Williams), and the native \te{చేయు}, \textquotedblleft{}to do\textquotedblright{} (DEDR 1957).
\end{cousinweb}

\begin{grammarlens}[title={What you've built}]
One more word for getting around.
\end{grammarlens}

\subsection*{Your turn}
\begin{itemize}
\raggedright
  \item \emph{Say it:} \emph{kaluṣitaṁ cēyu}
  \item \emph{Say it:} \emph{kaluṣitaṁ cēyu}, once more
  \item \emph{Say it:} \emph{karmāgārālu nadi nīṭini kaluṣitaṁ cēśāyi.} — The factories polluted the river water.
  \item \emph{Recall:} say the Telugu for waste, then the Telugu for a gas, then say \emph{kaluṣitaṁ cēyu} again
\end{itemize}

\begin{tcolorbox}[breakable,colback=teal!4,colframe=teal!35!black,title={Before you move on}]
What does \te{కలుషితం} \te{చేయు} mean? (\textquotedblleft{}To pollute\textquotedblright{}.) Say it once more.
\end{tcolorbox}
\section[\texorpdfstring{\te{స్వచ్ఛమైన} (svacchamaina)}{స్వచ్ఛమైన (svacchamaina)}]{\te{స్వచ్ఛమైన} (svacchamaina) — pure; clean}

\label{lesson:TE-C470-svacchamaina}



\begin{quote}
Before the new one: say the Telugu for a gas, then the Telugu for to pollute.
\end{quote}

\subsection*{You'll want to know: \te{స్వచ్ఛమైన}}
\textbf{\te{స్వచ్ఛమైన}} — \emph{svacchamaina} — \textquotedblleft{}pure; clean\textquotedblright{}.

\begin{quote}
\te{కొండల్లో} \te{స్వచ్ఛమైన} \te{గాలి} \te{దొరుకుతుంది}.
\end{quote}

\begin{quote}
\emph{koṇḍallō svacchamaina gāli dorukutundi.}
\end{quote}

\begin{quote}
In the hills you get pure air.
\end{quote}

It stands before its noun: \te{స్వచ్ఛమైన} \te{నీరు}, clean water.

\begin{cousinweb}
\textbf{\te{స్వచ్ఛమైన}} (\emph{svacchamaina}) is built from Sanskrit \emph{svaccha}, \textquotedblleft{}very clear, transparent\textquotedblright{} (Monier-Williams), with the Telugu adjective ending -\te{మైన}.
\end{cousinweb}

\begin{grammarlens}[title={What you've built}]
One more word for getting around.
\end{grammarlens}

\subsection*{Your turn}
\begin{itemize}
\raggedright
  \item \emph{Say it:} \emph{svacchamaina}
  \item \emph{Say it:} \emph{svacchamaina}, once more
  \item \emph{Say it:} \emph{koṇḍallō svacchamaina gāli dorukutundi.} — In the hills you get pure air.
  \item \emph{Recall:} say the Telugu for a gas, then the Telugu for to pollute, then say \emph{svacchamaina} again
\end{itemize}

\begin{tcolorbox}[breakable,colback=teal!4,colframe=teal!35!black,title={Before you move on}]
What does \te{స్వచ్ఛమైన} mean? (\textquotedblleft{}Pure; clean\textquotedblright{}.) Say it once more.
\end{tcolorbox}
